\section{A frozen continuation comparison under recursive risk}\label{sec:study95}
We apply the preceding construction to $d=2$ at both declared risk levels. There are four arms: a fitted oblique ReLU continuation, an additive linear-spline continuation, a projected-quadratic continuation, and conventional online evaluation of the original one-date risk Bellman problem. The last arm pays for that calculation; it is not supplied with a free optimal-value oracle. Every fitted label is itself computed from the original terminal primitive and a paid all-action certificate. The terminal continuation is the known $g$, so this last-date fit is consistent with the own-future construction rather than training against an externally supplied optimal continuation.

\subsection{Frozen design and common accuracy}
Each risk level has 1,024 training states, 128 calibration states and 2,048 evaluation states. The three streams are domain-separated before execution. States are drawn from the declared uniform 20-bit dyadic initial-state distribution; the state dynamics and innovations between decisions remain continuous. Each fitted arm uses the same labels, 32 fixed nonlinear features, an unrestricted affine part and nonnegative nonlinear coefficients. Least squares uses the same solver tolerance and iteration cap; readouts are stored on a 30-bit dyadic grid. The source, serialized coefficients and calibration selector are checked before evaluation. No fitting or reselection occurs on the evaluation stream.

The selector minimizes calibration mean cost-interval midpoints with a fixed lexicographic tie rule. It chooses online risk-DP at both risk levels. All twelve pairwise comparisons across the two levels were specified in advance, including those not involving the selected arm. The independent observation is an initial state, not a shock bin or a repeated solver query.

A common own-continuation gap is insufficient for a fair policy comparison. We require both an own-continuation gap at most $1/64$ and an original-optimum complete-policy loss bound at most $1/16$ at every queried initial state. Fitted arms pay for the extra original-continuation lower witness in Theorem~\ref{thm:continuous95}; online risk-DP reuses its already computed witness. Integration may refine from 128 to 4,096 bins, but the proposed first action cannot change. All failures are retained rather than silently replaced by conventional actions. The terminal controller is common to every arm and has the uniform allowance proved above.

\begin{table}[htbp]
\caption{Fitted continuations and whole-square residual enclosures}\label{tab:fits95}
\centering\small
\begin{tabular}{@{}llrrl@{}}
\toprule
$\alpha$ & Continuation & Rank & Training RMSE & $V_1-f$ on the whole square\\\midrule
0.5 & Quadratic & 6 & 0.04122 & [-0.13508, 0.52318]\\
0.5 & ReLU & 35 & 0.01668 & [-0.04373, 0.08665]\\
0.5 & Spline & 35 & 0.03979 & [-0.11690, 0.50298]\\
0.75 & Quadratic & 6 & 0.04166 & [-0.11726, 0.53606]\\
0.75 & ReLU & 35 & 0.01670 & [-0.04599, 0.09043]\\
0.75 & Spline & 35 & 0.03984 & [-0.10431, 0.45189]\\
\bottomrule
\end{tabular}
\par\smallskip\begin{minipage}{.98\linewidth}\footnotesize Each fit uses the same 1,024 labels and 32 nonlinear features plus an affine part. The quadratic feature matrix has rank six, not 35; the experiment matches nominal features and labels, not independent parameter rank. The residual enclosures cover every state in the continuous square.\end{minipage}
\end{table}

The ReLU fits have the smallest training RMSE and the narrowest whole-square residual brackets in both cases. These are distinct observations: training errors alone do not give a whole-domain certificate. The quadratic design has effective rank six although all three fitted arms have 35 stored coefficients; matching feature counts is therefore not matching approximation spaces or statistical degrees of freedom. No rank-matched representation superiority is inferred from this comparison.

\subsection{A common-policy comparison of certificate accounts}
Let $[l,u]$ be the verified whole-square bracket for $V_1-f$, $E=\max\{|l|,|u|\}$, and $\delta(x)$ the own-continuation all-action gap. Two valid original-optimum bounds for the same complete policy are
\[
 B_{\mathrm{norm}}(x)=\delta(x)+\beta(2E+\varepsilon_1),\qquad
 B_{\mathrm{signed}}(x)=\delta(x)+\beta(u-l+\varepsilon_1).
\]
Monotonicity and cash invariance prove these inequalities directly. The joint local-support bound is the paid bound in \eqref{eq:originalgap95}, not a new policy. Table~\ref{tab:accounts95} re-accounts for every existing fitted-arm row under all three bounds. At the fixed whole-square cover, neither scalar account meets $1/16$ on any of these rows, whereas the joint certificate meets it on all 12,288 fitted-arm rows. Thus keeping the economically relevant local continuation information matters for qualification, even though it does not establish a neural policy-cost advantage.

\begin{table}[htbp]
\caption{Three valid certificate accounts for the same frozen fitted policies}\label{tab:accounts95}
\centering\small
\begin{tabular}{@{}llrrrr@{}}\toprule
$\alpha$ & Continuation & Supremum norm & Signed scalar & Joint support & Qualified\\\midrule
0.5 & Quadratic & 0.981430 & 0.617589 & 0.005501 & 2048/2,048\\
0.5 & ReLU & 0.162986 & 0.122750 & 0.021836 & 2048/2,048\\
0.5 & Spline & 0.943607 & 0.581660 & 0.012112 & 2048/2,048\\
0.75 & Quadratic & 1.005610 & 0.612983 & 0.006422 & 2048/2,048\\
0.75 & ReLU & 0.170090 & 0.128429 & 0.022145 & 2048/2,048\\
0.75 & Spline & 0.847773 & 0.521923 & 0.008098 & 2048/2,048\\
\bottomrule\end{tabular}
\par\smallskip\begin{minipage}{.98\linewidth}\footnotesize
Entries are maximum original-optimum policy-loss bounds across the same 2,048 rows. Supremum-norm and signed-scalar bounds use the same frozen whole-square residual cover. Joint support uses the paid original-value witness. No row satisfies $1/16$ with either scalar account; every row does with the joint certificate. This is post-evaluation deterministic re-accounting, not a prospective runtime ablation or an optimality claim for any certificate-construction algorithm.
\end{minipage}\end{table}

This re-accounting was performed after the primary evaluation, with unchanged data, actors and targets. It is a deterministic certificate comparison, not a prospectively registered comparison of optimally tuned centered and uncentered solvers. The scalar accounts pay for the global residual cover; the joint account pays for the original-value witness. Their measured preparation and query work remain recorded separately, so certificate width is not converted into an assumed runtime saving.

\subsection{Original policy loss and economic accuracy}
All 16,384 arm-state records satisfy both qualification targets. The largest original-optimum loss bound among them is less than $0.022145$, below the common $0.0625$ target. This is a statement about every recorded initial state, not an empirical proof of that same tolerance on all possible initial states. The separate residual cover qualifies the whole continuous state square.

\begin{table}[htbp]
\caption{Original nested-risk cost and qualification at 2,048 evaluation states per risk level}\label{tab:cost95}
\centering\small
\begin{tabular}{@{}lllrr@{}}
\toprule
$\alpha$ & Continuation & Mean cost enclosure & Maximum loss bound & Failures\\\midrule
0.5 & Online risk-DP & [0.625634, 0.632502] & 0.003435 & 0\\
0.5 & Quadratic & [0.625754, 0.632622] & 0.005501 & 0\\
0.5 & ReLU & [0.626391, 0.633259] & 0.021836 & 0\\
0.5 & Spline & [0.625792, 0.632660] & 0.012112 & 0\\
0.75 & Online risk-DP & [0.625663, 0.632531] & 0.003435 & 0\\
0.75 & Quadratic & [0.625808, 0.632676] & 0.006422 & 0\\
0.75 & ReLU & [0.626600, 0.633468] & 0.022145 & 0\\
0.75 & Spline & [0.625805, 0.632674] & 0.008098 & 0\\
\bottomrule
\end{tabular}
\par\smallskip\begin{minipage}{.98\linewidth}\footnotesize Cost enclosures here contain numerical uncertainty only. Simultaneous sampling intervals for all paired differences appear in Table~\ref{tab:pairs95}. The maximum loss bound is against the original Bellman optimum; all displayed endpoints are rounded outwards. A failure is a failed own-continuation or original-optimum target, with every row retained.\end{minipage}
\end{table}

The estimand is the initial-distribution mean of each complete policy's original nested-risk value. It is not average realized cash expenditure or a monetary welfare estimate. The fixed margin $1/16$ equals four fifths of the maximal first-date investment charge $h(1/4)=5/64$ in this normalized model. It is a declared numerical economic tolerance; judging it small for a particular application requires an external economic calibration.

All twelve simultaneous empirical-Bernstein intervals lie strictly inside $[-1/16,1/16]$. Thus all four complete policies are equivalent at that prespecified tolerance on the declared initial distribution. None of the intervals excludes zero, so none establishes strict superiority. In particular, the narrower ReLU residual bracket does not translate here into an identified improvement over online risk-DP. The numerical mean-cost brackets in Table~\ref{tab:cost95} are not themselves sampling confidence intervals; the full paired intervals and their sampling components are reported in the appendix.

The numerical component includes original-law integration, terminal observation and optimization, and outward arithmetic. Initial-state discretization is part of the declared evaluation distribution, not an unreported approximation to a continuous uniform initial distribution. There is no transition-law substitution by a state grid, and therefore no hidden state-aggregation allowance in these costs. Sampling and numerical uncertainty are combined as proved in the supplement. The protocol does not permit changing the margin or sample size after examining the evaluation outcomes.

\subsection{Verification work and the value of storage}
\begin{table}[htbp]
\caption{Measured work of the completed evaluation attempt}\label{tab:work95}
\centering\small
\begin{tabular}{@{}llrrr@{}}
\toprule
$\alpha$ & Continuation & Seconds & Terminal proposals & All-action witnesses\\\midrule
0.5 & Online risk-DP & 32.509 & 112,721,920 & 2,048\\
0.5 & Quadratic & 12.254 & 13,631,488 & 4,096\\
0.5 & ReLU & 11.797 & 13,631,488 & 4,096\\
0.5 & Spline & 11.442 & 13,631,488 & 4,096\\
0.75 & Online risk-DP & 30.770 & 112,721,920 & 2,048\\
0.75 & Quadratic & 11.746 & 13,631,488 & 4,096\\
0.75 & ReLU & 11.237 & 13,631,488 & 4,096\\
0.75 & Spline & 11.523 & 13,631,488 & 4,096\\
\bottomrule
\end{tabular}
\par\smallskip\begin{minipage}{.98\linewidth}\footnotesize Every arm processes all 2,048 initial states. Timings include first-action proposals, qualification, the original-value witness, complete-policy cost enclosures and CSV output. They exclude preparation, the independent audit, diagnostic screening and failed earlier execution starts, all reported separately. Counts are instrumented work units, not every arithmetic instruction. This is a single x86-64 execution, not a cross-architecture speed guarantee.\end{minipage}
\end{table}

In the completed execution, the fitted arms require approximately 11.2--12.3 seconds per 2,048-state risk-level workload, compared with approximately 30.8--32.5 seconds for online risk-DP. These measurements include the extra original-optimum witness required of the fitted arms. Preparation for the entire two-risk-level, six-fit comparison takes approximately 8.79 seconds. A separately coded rational audit takes approximately 89.98 seconds and checks the continuous-risk kernels and retained evidence; that verification is a separate paid campaign cost, not zero-cost validation.

These measurements identify a benefit of stored continuations on this workload, not a uniquely neural advantage. Spline and quadratic continuations obtain similar execution reductions, and the fitted arms are not rank matched. The counts identify the repeated terminal computations avoided by storage; they do not identify the cost of every machine instruction. The reported timings are from one x86-64 execution. Two earlier tool-interrupted starts are recorded separately, without recoverable complete cost ledgers. Accordingly the table is not a claim about the aggregate cost of every attempted run, nor a cross-architecture, energy or asymptotic scalability result.

Whole-bin screening excludes between 738 and 746 of 4,096 candidates at $\alpha=1/2$, and between 2,053 and 2,062 at $\alpha=3/4$. After paying for unpruned original-value bounds and the strict test, the screened calculation is slightly slower in every displayed case. The exact unchanged tail values validate the exclusions; they do not establish a computational gain. No-cut rows and all denominators are retained in the appendix. The experiment therefore separates a valid geometric certificate, the amount it removes and the net value of constructing it.

The two-date scalar-control experiment does not replace the inherited two-control, two-innovation horizon frontier. Proposition~\ref{prop:horizon95} now gives explicit fixed-accuracy resolutions and a horizon-uniform construction for both the scalar and original vector economies. The general NBO construction and joint-domain theorem remain available under their hypotheses, while a full independent largest-vector implementation remains a separate numerical obligation. This distinction preserves the general research program without borrowing scalability evidence from a different problem.
